% Optional manuscript fragment. All formulas below are existing conjectures
% in the pinned source; none is asserted to be proved by the present report.
\subsection{The surviving named Gaussian conjectures}
\label{sec:surviving-gaussian-conjectures}

The manuscript's names $S_6$ and $S_8$ refer to specific alternating
harmonic sums, and should not be confused with the uncolored sums of our
Gamma-ratio theorem.  With the largest summation index first in the
multiple-polylogarithm convention, their definitions are
\[
 S_p=\sum_{n\geq0}\frac{(-1)^nH_n}{(2n+1)^p},\qquad
 g_{a,b}=\operatorname{Im}\operatorname{Li}_{a,b}(i,1)
 =\sum_{n\geq0}\frac{(-1)^nH_{2n}^{(b)}}{(2n+1)^a},
\]
where $H_m^{(b)}=\sum_{j=1}^{m}j^{-b}$ and $H_m=H_m^{(1)}$.
Write $G=\beta(2)$ for Catalan's constant.  The existing mixed-endpoint
identity gives
\[
 S_p=g_{p,1}+\operatorname{Im}\operatorname{Li}_{p,1}(i,-1).
\]
In particular, $S_p$ is not just $g_{p,1}$.

The frozen weight-seven conjecture currently reads
\begin{align}
 S_6\stackrel{?}{={}}&\frac{722}{527}g_{6,1}
 +\frac{40}{527}g_{4,3}-\frac{128}{155}g_{2,5}
 +\frac{15191}{28569600}\pi^7\notag\\
 &-\frac3{124}G\zeta(5)
 -\frac{2373}{10540}\beta(4)\zeta(3)-2\beta(6)\log2.
 \label{eq:source-current-S6}
\end{align}
The source is \nolinkurl{chapters/04-cyclotomic-quotients.tex},
lines 338--349.  Its label is \texttt{cycloquot:conj:S6}.
The alternative formula involving $g_{5,2}$ is already proved equivalent
to this one by two convergent shuffle identities in the same file,
lines 295--336.  That coordinate change does not prove either conjectured
vanishing.  Moreover, the source's separating functional is nonzero on
the frozen relation, so a proof cannot use only the restricted
single-product presentation considered there.

The current weight-nine conjecture is
\begin{align}
 S_8\stackrel{?}{={}}&\frac{52334}{37973}g_{8,1}
 +\frac{2824}{37973}g_{6,3}
 -\frac{6832}{189865}g_{4,5}
 -\frac{139760}{265811}g_{2,7}\notag\\
 &+\frac{998183\pi^9}{17283022848}
 -\frac{78615}{19442176}G\zeta(7)
 -\frac{201021}{4252976}\beta(4)\zeta(5)\notag\\
 &-\frac{290505}{1063244}\beta(6)\zeta(3)-2\beta(8)\log2.
 \label{eq:source-current-S8}
\end{align}
Its exact source is \nolinkurl{chapters/04-S8-candidate.tex},
lines 7--17, label \texttt{s8new:conj:S8}.
The same file, lines 40--50, records exact rational proximity bounds
below $10^{-355}$ for the normalized residuals of both current
conjectures.  Both enclosures contain zero.  The equations remain
conjectural in the pinned manuscript and in the present report.

This $S_8$ candidate must be distinguished from the earlier frozen vector
in the basket containing $g_{7,2}$ and $g_{5,4}$.  That different vector
has a strictly negative certified residual between
$-1.3274179020580459\times10^{-86}$ and
$-1.3274179020580457\times10^{-86}$; see
\nolinkurl{chapters/10-discovery.tex}, lines 259--301,
label \texttt{research:prop:S8-rejected}.
The rejection concerns precisely that vector and gives no arithmetic
independence result.

Our uncolored harmonic theorem does not acquire the factor $(-1)^n$
merely by setting its Hurwitz parameter to $1/2$.  Its direct application
therefore cannot prove either displayed Gaussian identity.  The
trilinear theorem does provide root-of-unity colored Tornheim
coordinates at rational shifts, with multiple-polylogarithm reductions
at positive integer inner orders.  A coordinate representation alone
does not annihilate either frozen residual, and the partial-fraction
identities for discrete inner orders cannot be differentiated as though
they held with a continuously varying summation range.  The remaining
concrete problem is to find and prove the missing mixed-endpoint
functional relation at weights seven and nine, and then reduce it
exactly to the specified frozen basket.
